\subsection{The volume detected by ideals}
\label{sec:ch51-current-volume}

\begin{definition}[Volume detected by ideals]
\label{def:complex-current-volume}
For \(\varphi\in\PSH(X,\theta)\), define
\[
 \vol(\theta,\varphi)
 :=\int_X\bigl(\theta+\ddc
             P_\theta[\varphi]_{\mathcal I}\bigr)^n.
\tag{51.32}\label{eq:complex-current-volume}
\]
It depends only on the current \(\theta_\varphi\).  Moreover,
\[
 M_\theta(\varphi)\leq\vol(\theta,\varphi),
\tag{51.33}\label{eq:complex-mass-volume}
\]
with equality for analytic singularities.
\end{definition}

Three comments.  The nonpluripolar product on the right of
\eqref{eq:complex-current-volume} is, as always, the one of
\eqref{eq:complex-global-npp}; the notation omits the angle brackets only
because the argument is now a specific model potential.  The inequality
\eqref{eq:complex-mass-volume} is
\Cref{thm:complex-npp-ledger} together with
\(P_\theta[\varphi]\leq P_\theta[\varphi]_{\mathcal I}\) from
\eqref{eq:complex-p-i-comparison} and the mass identity
\eqref{eq:complex-p-projection}; explicitly,
\[
 M_\theta(\varphi)
 =M_\theta\bigl(P_\theta[\varphi]\bigr)
 \leq M_\theta\bigl(P_\theta[\varphi]_{\mathcal I}\bigr)
 =\vol(\theta,\varphi),
\]
the middle step because \(P_\theta[\varphi]\preceq
P_\theta[\varphi]_{\mathcal I}\).  Equality for analytic singularities is
\eqref{eq:complex-analytic-envelope-agreement}.  Finally, the dependence
only on the current is the statement that
\(\vol(\theta,\varphi)\) is unchanged when \(\varphi\) is replaced by
\(\varphi+c\), which is
\Cref{prop:ch51-mass-inequalities}, item
\cref{itm:ch51-mi:translation}, together with the fact that
\eqref{eq:complex-i-envelope} depends only on the
\(\mathcal I\)-class, which is invariant under adding constants.
